\chapter{Enchères}\label{ch:m1:auctions}

À seize heures moins dix à New York, le carnet d'ordres d'une grande action ressemble à ce qu'il
a été toute la journée~: quelques centaines d'actions à chaque cours, à un cent d'écart. À côté,
invisibles sur l'écran ordinaire, des ordres d'achat et de vente portant sur plusieurs millions
d'actions se sont accumulés avec une seule consigne~: \emph{au \omterm{def:m1:auctions:moc}{cours de clôture}, quel qu'il soit}. À seize
heures, tous seront exécutés, à un seul prix, au même instant, pour une quantité que le marché continu
n'aurait pas pu absorber en une heure. Ce prix évaluera chaque fonds qui détient le titre, réglera
les options qui arrivent à échéance aujourd'hui et servira de référence aux limites du lendemain.
Il est calculé par un algorithme d'une trentaine de lignes.
Ce chapitre présente cet algorithme, l'information publiée pendant son attente,
et les raisons pour lesquelles une seule cotation par jour
est devenue la plus importante.

\section{L'enchère à la criée}

\begin{definition}[Enchère et prix d'équilibre]\label{def:m1:auctions:call}
Dans une \emph{enchère}\index{enchère} les ordres sont recueillis pendant une période
sans être exécutés~; à la fin de celle-ci, un seul
\emph{prix d'équilibre}\index{prix d'équilibre} est calculé, et chaque ordre pouvant être exécuté
à ce prix l'est, à ce prix. Les enchères d'ouverture et de clôture, les \omterm{def:m1:auctions:volint}{interruptions de volatilité}
et les introductions en bourse sont des enchères~; de même, en miniature,
les \omterm{def:m1:european-equity-market-structure:periodic}{enchères périodiques} de la
\cref{def:m1:european-equity-market-structure:periodic}.
\end{definition}

Pour un prix candidat $p$, notons $D(p)$ la quantité totale d'ordres d'achat
prêts à payer $p$ ou davantage (ordres au marché compris) et $S(p)$ la quantité totale
d'ordres de vente prêts à accepter $p$ ou moins. $D$ est décroissante, $S$ croissante~;
la quantité pouvant être échangée à $p$ est $\min(D(p), S(p))$.

\begin{method}[Les règles de détermination du prix d'équilibre]\label{met:m1:auctions:rules}
Parmi les cours limites présents dans le carnet et le cours de référence~:
\begin{enumerate}
\item \textbf{Volume.} Retenir les prix qui maximisent le volume exécutable
  $\min(D(p),S(p))$.
\item \textbf{Surplus.} Parmi ceux-ci, retenir les prix qui minimisent la quantité non exécutée
  $|D(p) - S(p)|$.
\item \textbf{Pression du marché.} Si à chaque prix restant le surplus est du côté acheteur,
  retenir le plus élevé~; s'il est du côté vendeur, retenir le plus bas.
\item \textbf{Cours de référence.} Sinon, retenir le prix le plus proche du cours de référence
  (généralement le dernier cours échangé).
\end{enumerate}
Puis répartir~: les ordres au marché d'abord, puis les ordres à cours limité
meilleurs que le \omterm{def:m1:auctions:call}{prix d'équilibre}, puis les ordres au \omterm{def:m1:auctions:call}{prix d'équilibre} dans l'ordre de priorité
temporelle, jusqu'à épuisement du volume exécutable de chaque côté.
\end{method}

Les deux premières règles sont communes aux modèles de marché publiés par les bourses
européennes~: le prix d'\omterm{def:m1:auctions:call}{enchère} est celui qui offre «~le plus de volume exécutable et le plus faible
surplus~». Les règles de départage varient dans le détail d'une plateforme à l'autre, et une
implémentation doit suivre le règlement de marché de la plateforme qu'elle simule~; la logique ci-dessus
est la plus courante.

\begin{example}[Huit ordres]\label{ex:m1:auctions:eight}
Acheteurs~: 300 au marché~; 500 à 10.03~; 400 à 10.01~; 600 à 10.00. Vendeurs~: 200
au marché~; 400 à 9.99~; 500 à 10.01~; 700 à 10.02. Le cours de référence est
10.00.
\begin{center}\small
\begin{tabular}{@{}lrrrr@{}}
\hline
Prix & $D(p)$ & $S(p)$ & Exécutable & Surplus \\
\hline
9.99 & 1\,800 & 600 & 600 & $+1\,200$ \\
10.00 & 1\,800 & 600 & 600 & $+1\,200$ \\
10.01 & 1\,200 & 1\,100 & \textbf{1\,100} & $+100$ \\
10.02 & 800 & 1\,800 & 800 & $-1\,000$ \\
10.03 & 800 & 1\,800 & 800 & $-1\,000$ \\
\hline
\end{tabular}
\end{center}
La règle 1 seule retient 10.01~: 1\,100 actions sont échangées. Du côté acheteur, l'ordre au marché
(300) et l'ordre à 10.03 (500) sont exécutés en totalité, et l'ordre à 10.01 reçoit 300 de ses 400~;
du côté vendeur, tout ce qui est jusqu'à 10.01 est exécuté.
Un surplus acheteur de 100 actions subsiste à 10.01, et c'est ce que la bourse publie
comme déséquilibre de clôture.
\end{example}

\begin{omfigure}
\begin{tikzpicture}
\begin{axis}[width=10.5cm,height=5.2cm,
  xlabel={prix},ylabel={quantité (actions)},
  tick label style={font=\small},label style={font=\small},
  xmin=9.965,xmax=10.055,ymin=0,ymax=2000,grid=major,grid style={gray!25},
  xtick={9.97,9.98,9.99,10.00,10.01,10.02,10.03,10.04,10.05},
  xticklabel style={/pgf/number format/fixed,/pgf/number format/precision=2,/pgf/number format/fixed zerofill},
  yticklabel style={/pgf/number format/1000 sep={\,}},
  legend columns=2,legend style={font=\footnotesize,at={(0.5,-0.3)},anchor=north,draw=none,fill=none,/tikz/every even column/.append style={column sep=8pt}},
  table/col sep=comma,x filter/.code={\pgfmathparse{\pgfmathresult/100}}]
\addplot[omDef,very thick,const plot mark mid,mark=*,mark size=1.3pt] table[x=price,y=demand]{figdata/markets-1/13-auctions/curves.csv};
\addplot[omThm,very thick,const plot mark mid,mark=*,mark size=1.3pt] table[x=price,y=supply]{figdata/markets-1/13-auctions/curves.csv};
\addplot[only marks,mark=o,mark size=4pt,black,thick] coordinates {(1001,1100)};
\legend{demande $D(p)$,offre $S(p)$}
\end{axis}
\end{tikzpicture}

\omcaption{Demande et offre de l'\cref{ex:m1:auctions:eight}. Les courbes sont des fonctions
en escalier du prix~; l'\omterm{def:m1:auctions:call}{enchère} échange la plus grande quantité sur laquelle elles peuvent
s'accorder, soit 1\,100 actions à 10.01 (cerclé~: c'est là que la plus basse des deux
courbes est la plus haute). Données~: le carnet d'exemple
du chapitre.}\label{fig:m1:auctions:curves}
\end{omfigure}

\begin{proposition}[Pourquoi un prix unique]\label{prop:m1:auctions:welfare}
Parmi tous les ensembles de transactions entre les ordres du carnet
qui respectent la limite de chaque ordre, l'\omterm{def:m1:auctions:call}{enchère} à un prix satisfaisant la règle 1
maximise le volume échangé, et chaque participant échange à un prix au moins aussi bon
que sa limite, sans qu'aucun participant puisse se plaindre qu'un ordre moins bien coté
ait été exécuté à sa place.
\end{proposition}
\begin{proof}
Tout ensemble réalisable de transactions à des prix quelconques associe un acheteur
de limite $b$ à un vendeur de limite $s \le b$. Ordonnons les acheteurs par limite décroissante
et les vendeurs par limite croissante~: la $k$-ième paire est réalisable si et seulement si
le $k$-ième meilleur cours acheteur est au moins égal au $k$-ième meilleur cours vendeur,
et le plus grand tel $k$ (en actions) vaut $\max_p \min(D(p),S(p))$, atteint en appariant
dans cet ordre. Échanger toutes ces paires à un seul prix $p^\star$ compris entre le cours acheteur
marginal et le cours vendeur marginal respecte chaque limite, et tout ordre exclu a une limite
non meilleure que celle d'un ordre inclus.
\end{proof}

\section{Ouverture, clôture et le reste}

\begin{definition}[Cours de clôture, ordres au marché à la clôture et ordres à cours limité à la clôture]\label{def:m1:auctions:moc}
Le \emph{cours de clôture}\index{cours de clôture} d'un titre est le prix de
son \omterm{def:m1:auctions:call}{enchère} de clôture sur la bourse où il est coté. Un
\emph{ordre au marché à la clôture}\index{ordre au marché à la clôture} (MOC) s'exécute dans
cette \omterm{def:m1:auctions:call}{enchère}, à quelque prix qu'elle fixe~; un
\emph{ordre à cours limité à la clôture}\index{ordre à cours limité à la clôture} (LOC) s'exécute dans cette \omterm{def:m1:auctions:call}{enchère}
seulement à son cours limite ou à un cours meilleur.
\end{definition}

\begin{definition}[Déséquilibre des ordres et cours indicatif]\label{def:m1:auctions:imbalance}
Pendant la période d'\omterm{def:m1:auctions:call}{enchère}, la bourse publie le
\emph{cours indicatif}\index{cours indicatif}, soit le \omterm{def:m1:auctions:call}{prix d'équilibre}
auquel l'\omterm{def:m1:auctions:call}{enchère} serait fixée si elle prenait fin maintenant, la quantité \emph{appariée} à ce
prix, et le \emph{déséquilibre des ordres}\index{déséquilibre des ordres}~: le sens et
la taille de la quantité qui n'a pas pu être appariée.
\end{definition}

\begin{dated}{2026-09}{Les dix dernières minutes à New York}\label{dat:m1:auctions:us}
Au New York Stock Exchange, les ordres MOC et LOC peuvent être saisis
jusqu'à 15:50~; à cette heure, la bourse publie un déséquilibre réglementaire pour chaque
valeur dont le déséquilibre atteint au moins 500 \omterm{def:m1:us-equity-market-structure:lots}{lots standard}, après quoi les ordres de clôture
ne sont acceptés que du côté opposé, et l'information sur le déséquilibre est
republiée chaque seconde. Sur le Nasdaq, l'information sur le déséquilibre est diffusée
à partir de 15:50, toutes les cinq secondes puis chaque seconde à partir de 15:55~; les ordres MOC
sont acceptés jusqu'à 15:55 et les ordres LOC jusqu'à 15:58. Les deux \omterm{def:m1:auctions:call}{enchères} se déroulent à
16:00.
\end{dated}

\begin{omfigure}
\begin{tikzpicture}[font=\small]
\draw[-{Stealth[length=2.2mm]},thick,omExo] (-0.4,0) -- (15.2,0);
\foreach \x/\t in {1.0/15:50,6.0/15:55,10.0/15:58,14.0/16:00}{
  \draw[omExo,thick] (\x,-0.12) -- (\x,0.12);
  \node[font=\small\bfseries,anchor=north] at (\x,-0.2) {\t};
}
\node[omDef,font=\footnotesize,align=center,anchor=south] at (1.0,0.25) {publication du\\déséquilibre~; le NYSE\\clôt les MOC et LOC};
\node[omDef,font=\footnotesize,align=center,anchor=south] at (6.0,0.25) {Nasdaq clôt\\la saisie des MOC};
\node[omDef,font=\footnotesize,align=center,anchor=south] at (10.0,0.25) {Nasdaq clôt\\la saisie des LOC};
\node[omThm,font=\footnotesize,align=center,anchor=south] at (14.0,0.25) {les deux enchères\\fixent le prix d'équilibre};
\draw[omThm,thick,decorate,decoration={brace,mirror,amplitude=5pt}] (1.0,-0.95) -- (14.0,-0.95)
  node[midway,below=6pt,font=\footnotesize,omThm] {le déséquilibre est publié~; le marché continu se poursuit};
\end{tikzpicture}

\omcaption{Les dix dernières minutes de séance sur les deux bourses américaines de cotation, d'après
l'\cref{dat:m1:auctions:us}. Les heures limites existent pour que le déséquilibre publié
ait un sens~: une fois l'heure limite d'une bourse passée, les nouveaux ordres de clôture
qui y sont saisis ne peuvent que la réduire.}\label{fig:m1:auctions:timeline}
\end{omfigure}

L'\emph{\omterm{def:m1:auctions:call}{enchère} d'ouverture} répond à un problème différent~: les nouvelles survenues pendant la nuit ont
déplacé le prix juste d'une quantité inconnue, et les premières négociations en continu
se feraient à l'aveugle. Une période d'\omterm{def:m1:auctions:call}{enchère} permet à chacun de voir le \omterm{def:m1:auctions:imbalance}{cours indicatif}
converger avant de s'engager. Elle attire bien moins de volume que la clôture,
car personne n'a l'ouverture pour indice de référence.

\begin{definition}[Interruption de volatilité]\label{def:m1:auctions:volint}
Une \emph{interruption de volatilité}\index{interruption de volatilité} est une
\omterm{def:m1:auctions:call}{enchère} non programmée qu'une bourse déclenche automatiquement lorsque la transaction suivante
se produirait en dehors d'un couloir de prix autour d'un prix de référence~: la
négociation en continu s'arrête pendant quelques minutes, les ordres s'accumulent et
la négociation reprend au \omterm{def:m1:auctions:call}{prix d'équilibre}.
\end{definition}

La réponse des bourses européennes à un mouvement brutal consiste donc à modifier le mécanisme
plutôt que d'interdire les prix, comme le fait une limite quotidienne
(la \cref{def:m1:asian-and-emerging-equity-markets:limit}), ou à suspendre la
négociation, comme le font les bandes Limit Up--Limit Down (LULD) (au \cref{ch:m1:stress-case-studies}).

\section{Les déséquilibres sont de l'information}

\begin{omfigure}
\begin{tikzpicture}
\begin{axis}[width=10.5cm,height=5.2cm,
  xlabel={déséquilibre acheteur ajouté (\% du volume équilibré)},ylabel={variation du cours de clôture (ticks)},
  tick label style={font=\small},label style={font=\small},
  xmin=0,xmax=100,ymin=0,ymax=4,grid=major,grid style={gray!25},table/col sep=comma]
\addplot[omDef,very thick,mark=*,mark size=1.4pt] table[x=imbalance_pct,y=move_ticks]{figdata/markets-1/13-auctions/impact.csv};
\end{axis}
\end{tikzpicture}

\omcaption{Effet d'un ordre d'achat à la clôture inattendu sur le cours, dans des
\omterm{def:m1:auctions:call}{enchères} simulées de 400 ordres à cours limité. La relation est presque linéaire
parce que la courbe d'offre au voisinage du \omterm{def:m1:auctions:call}{prix d'équilibre} l'est~; sa pente est la
profondeur de l'\omterm{def:m1:auctions:call}{enchère}. Données~: la simulation du tutoriel, 60 carnets par
point.}\label{fig:m1:auctions:impact}
\end{omfigure}

\begin{proposition}[Le prix d'un déséquilibre]\label{prop:m1:auctions:depth}
Supposons qu'au voisinage du \omterm{def:m1:auctions:call}{prix d'équilibre} le carnet présente une densité constante~: chaque
tick de prix ajoute $\lambda_S$ actions à l'offre et retire $\lambda_D$
actions de la demande. Un ordre d'achat au marché de $q$ actions ajouté à une
\omterm{def:m1:auctions:call}{enchère} équilibrée fait monter le \omterm{def:m1:auctions:call}{prix d'équilibre} d'environ $q/(\lambda_D +
\lambda_S)$ ticks.
\end{proposition}
\begin{proof}
Au prix initial $p_0$, $D = S$. En $p_0 + x$, la nouvelle demande vaut $D - \lambda_D
x + q$ et l'offre $S + \lambda_S x$~; elles sont égales en $x = q/(\lambda_D +
\lambda_S)$.
\end{proof}

Le dénominateur est la seule chose qui protège un \omterm{def:m1:auctions:moc}{cours de clôture} contre un gros
ordre, et il est endogène~: il se compose d'intervenants qui surveillent le
déséquilibre publié et introduisent un intérêt de sens contraire \emph{parce que} ils anticipent
que le déséquilibre fera bouger le prix. Les fournisseurs de liquidité à la clôture sont rémunérés par
l'écart entre le \omterm{def:m1:auctions:moc}{cours de clôture} et le cours quelques minutes avant
ou le lendemain matin~; les fonds indiciels le paient sciemment, en contrepartie
d'un \omterm{def:m1:the-buy-side:benchmark}{tracking error} nul (au \cref{ch:m1:index-construction-and-rebalancing}). Un ordre
conçu pour faire bouger la clôture plutôt que pour négocier --- «~marquer la clôture~» ---
est une manipulation de marché partout, et la transparence de l'\omterm{def:m1:auctions:call}{enchère} de clôture
est ce qui la rend détectable.

\section{Pourquoi la clôture continue de croître}

Trois forces poussent les volumes vers l'\omterm{def:m1:auctions:call}{enchère} de clôture, et chacune renforce
les deux autres. \emph{Indices de référence~:} les fonds indiciels, et tout gérant mesuré par rapport à un
indice, sont valorisés aux \omterm{def:m1:auctions:moc}{cours de clôture}~; négocier à un autre cours revient à prendre un
risque de \omterm{def:m1:the-buy-side:benchmark}{tracking error} sans aucune rémunération (au \cref{ch:m1:the-buy-side}).
\emph{Taille~:} l'\omterm{def:m1:auctions:call}{enchère} offre la liquidité la plus profonde de la journée à fourchette nulle,
si bien que les gros ordres indifférents au moment de l'exécution la préfèrent~; leur présence
la rend plus profonde encore. \emph{Amincissement ailleurs~:} le volume qui migre vers la
clôture laisse la séance continue moins profonde, ce qui renchérit le coût de
la négociation en cours de journée et envoie davantage d'ordres vers la clôture. En Europe, un quart du
volume négocié en bourse transite désormais par les \omterm{def:m1:auctions:call}{enchères} de clôture
(\cref{dat:m1:european-equity-market-structure:july}). La bourse de cotation possède cette \omterm{def:m1:auctions:call}{enchère} et, avec elle,
l'un des derniers leviers de pouvoir de fixation des prix dans la
négociation d'actions.

\section{Tutoriel~: calculer le prix d'équilibre d'un carnet}\label{tut:m1:auctions:uncross}

\begin{tutorial}
\textbf{Objectif.} Implémenter les règles de détermination du \omterm{def:m1:auctions:call}{prix d'équilibre}, reproduire
l'\cref{ex:m1:auctions:eight}, puis observer un \omterm{def:m1:auctions:imbalance}{cours indicatif} converger.
\textbf{État final~:} prix 10.01, volume 1\,100, surplus $+100$, et les deux
graphiques de cette section.
\begin{tutsteps}
\item \textbf{Les quatre règles}, sur des prix entiers exprimés en ticks.
\omcode{code/firm/auction/firm_auction.py}{38}{53}{Détermination du prix d'équilibre~: volume, surplus, pression du marché, prix de référence.}\label{lst:m1:auctions:uncross}
\item \textbf{Allocation.} Trier chaque côté de sorte que les ordres au marché viennent en premier,
puis les meilleurs prix, puis l'arrivée la plus ancienne, et exécuter jusqu'à épuisement
du volume.
\omcode{code/firm/auction/firm_auction.py}{56}{69}{Allocation avec priorité de prix et de temps.}\label{lst:m1:auctions:allocate}
\item \textbf{Exécuter l'exemple}~; les tests d'acceptation couvrent aussi l'absence
de croisement, une pression unilatérale, une égalité tranchée par le prix de référence, et
la priorité temporelle au \omterm{def:m1:auctions:call}{prix d'équilibre}.
\item \textbf{Une phase d'\omterm{def:m1:auctions:call}{enchère}.} Alimenter avec 600 ordres aléatoires en quarante paquets et
recalculer après chacun~: le \omterm{def:m1:auctions:imbalance}{cours indicatif} oscille de deux ou trois ticks au
départ, puis se stabilise à mesure que le volume apparié croît.
\end{tutsteps}
\textbf{Pour aller plus loin.} Ajouter une règle selon laquelle le \omterm{def:m1:auctions:call}{prix d'équilibre} doit se situer
dans une bande de prix autour de la dernière transaction en séance continue, et prolonger l'\omterm{def:m1:auctions:call}{enchère} sinon
--- une \omterm{def:m1:auctions:volint}{interruption de volatilité} au sein de l'\omterm{def:m1:auctions:call}{enchère}. Ajouter ensuite des ordres
de déséquilibre uniquement, qui ne peuvent s'exécuter que contre le surplus.
\end{tutorial}

\begin{omfigure}
\begin{tikzpicture}
\begin{axis}[width=10.5cm,height=4.8cm,
  xlabel={paquet d'ordres arrivés},ylabel={cours indicatif (ticks)},
  tick label style={font=\small},label style={font=\small},
  xmin=1,xmax=40,ymin=996,ymax=1004,grid=major,grid style={gray!25},
  yticklabel style={/pgf/number format/1000 sep={\,}},axis y line*=left,table/col sep=comma]
\addplot[omDef,very thick,const plot] table[x=step,y=price]{figdata/markets-1/13-auctions/indicative.csv};
\end{axis}
\begin{axis}[width=10.5cm,height=4.8cm,
  ylabel={volume apparié (milliers d'actions)},tick label style={font=\small},label style={font=\small},
  xmin=1,xmax=40,ymin=0,ymax=160,axis y line*=right,axis x line=none,table/col sep=comma]
\addplot[omThm,thick,dashed] table[x=step,y=volume_k]{figdata/markets-1/13-auctions/indicative.csv};
\end{axis}
\end{tikzpicture}

\omcaption{Une phase d'\omterm{def:m1:auctions:call}{enchère} simulée. Le \omterm{def:m1:auctions:imbalance}{cours indicatif} (en bleu, axe de gauche)
évolue librement tant que peu de volume est apparié, et se stabilise dès que le volume apparié
(en rouge pointillé, axe de droite) est grand par rapport à un ordre isolé quelconque. Données~: la
simulation du tutoriel.}\label{fig:m1:auctions:indicative}
\end{omfigure}

\section{Construction~: l'enchère}\label{bld:m1:auctions:uncross}

\begin{build}
\bfield{Rôle} Le simulateur de bourse du livre~10 ouvre et clôt
chaque séance par une \omterm{def:m1:auctions:call}{enchère}, et suspend la négociation continue par une
\omterm{def:m1:auctions:call}{enchère} lorsque les prix évoluent trop vite. Il s'agit de son moteur de détermination du \omterm{def:m1:auctions:call}{prix d'équilibre}. \bfield{Interface} \texttt{AuctionOrder(order\_id, side, quantity, price,
seq)} avec \texttt{price = None} pour un ordre au marché~;
\texttt{uncross(orders, reference)} renvoyant le prix, le volume, le surplus signé
et la liste des exécutions~; \texttt{demand} et \texttt{supply} pour publier
les informations indicatives.
\bfield{Règles} Prix entiers. Les quatre règles de la \cref{met:m1:auctions:rules}
dans cet ordre~; les prix candidats sont les limites présentes au carnet et le
prix de référence. Sans croisement, il n'y a ni prix ni exécution. Les quantités achetées et vendues
dans les exécutions sont toujours égales.
\bfield{Tests d'acceptation} \texttt{code/firm/auction/tests/}, \texttt{cpp/} et
\texttt{rust/}~: l'exemple du chapitre avec son exécution partielle~; absence de croisement~; pression acheteuse et
pression vendeuse~; égalité départagée par le prix de référence~; priorité temporelle. Les trois
implémentations doivent concorder.
\bfield{Pour aller plus loin} Ordres de déséquilibre uniquement et \omterm{def:m1:auctions:moc}{ordres à cours limité à la clôture} (LOC), assortis de
heures limites d'entrée, et une fonction qui publie la quantité appariée, le déséquilibre et
le \omterm{def:m1:auctions:imbalance}{cours indicatif} au format de l'\cref{dat:m1:auctions:us}.
\end{build}

\begin{omsources}
\item New York Stock Exchange, \emph{NYSE Closing Process} et \emph{NYSE
Opening and Closing Auctions} fiches d'information~; \emph{NYSE Closing Auction: Timing
Shifts and Marketability Trends}, 18 novembre 2025.
\item Nasdaq, \emph{The Nasdaq Opening and Closing Crosses: Frequently Asked
Questions}~; US Securities and Exchange Commission, Release 34-84454 (heures
limites d'entrée des ordres à la clôture).
\item Deutsche Börse, \emph{Xetra Market Model: Continuous Trading and
Auction}~; entrée de glossaire «~Auction principle~».
\item Euronext, \emph{Trading Manual for the Cash Market}.
\item A.~Madhavan, «~Trading mechanisms in securities markets~»,
\emph{Journal of Finance} 47 (1992).
\item M.~Pagano et R.~Schwartz, «~A closing call's impact on market
quality at Euronext Paris~», \emph{Journal of Financial Economics} 68
(2003).
\end{omsources}

\section{Exercices}

\begin{exercise}[$\star$]\label{exo:m1:auctions:1}
Ordres d'achat~: 400 à 20.10, 300 à 20.05, 500 à 20.00. Ordres de vente~: 200 à
19.95, 600 à 20.05, 300 à 20.10. Tabuler $D$, $S$ et le volume
exécutable à chaque prix limite, et donner le \omterm{def:m1:auctions:call}{prix d'équilibre} et le volume.
\end{exercise}

\begin{exercise}[$\star$]\label{exo:m1:auctions:2}
Pour l'\omterm{def:m1:auctions:call}{enchère} de l'exercice précédent, donner l'excédent et son côté, et
indiquer quels ordres sont exécutés, en totalité ou en partie.
\end{exercise}

\begin{exercise}[$\star$]\label{exo:m1:auctions:3}
En utilisant l'\cref{dat:m1:auctions:us}~: à 15:53, un trader veut ajouter un
ordre d'achat au marché à la clôture (MOC) sur une action cotée sur chacune des bourses. Sur laquelle peut-elle le faire,
et à quelle condition sur l'autre~?
\end{exercise}

\begin{exercise}[$\star\star$]\label{exo:m1:auctions:4}
Un carnet ne contient qu'un seul ordre d'achat, de 1\,000 actions à 50.20, et un seul
ordre de vente, de 1\,000 actions à 50.00. Le prix de référence est de 50.12. Appliquer les quatre
règles. Quel serait le prix avec une référence de 50.40~?
\end{exercise}

\begin{exercise}[$\star\star$]\label{exo:m1:auctions:5}
Au voisinage de son \omterm{def:m1:auctions:call}{prix d'équilibre}, une \omterm{def:m1:auctions:call}{enchère} présente 40\,000 actions d'offre et 25\,000
actions de demande par tick. Estimer l'effet d'un achat imprévu
au marché à la clôture de 300\,000 actions. Le tick vaut 1 centime et le titre
se négocie à 80~\$~: exprimer ce résultat en points de base.
\end{exercise}

\begin{exercise}[$\star\star$]\label{exo:m1:auctions:6}
Un fonds indiciel doit acheter pour 50 millions de dollars d'un titre à la clôture. Le
déséquilibre publié laisse penser que le \omterm{def:m1:auctions:moc}{cours de clôture} sera de 12 points de base au-dessus de
cours continu actuel~; acheter maintenant sur le marché continu coûterait un impact estimé
de 8 points de base. Quel est le coût espéré de chaque choix, et pourquoi
le fonds recourra-t-il quand même à la clôture~?
\end{exercise}

\begin{exercise}[$\star\star\star$]\label{exo:m1:auctions:7}
\emph{Codage.} Avec \texttt{imbalance\_impact}, estimer la pente du \omterm{def:m1:auctions:moc}{cours de clôture}
par rapport au déséquilibre pour des carnets de 400 et de 1\,600 ordres (60 carnets
chacun, déséquilibre égal à 50~\% du volume équilibré). Comparer les pentes
\emph{par action de déséquilibre} et les relier à la
\cref{prop:m1:auctions:depth}.
\end{exercise}

\begin{exercise}[$\star\star\star$]\label{exo:m1:auctions:8}
\emph{Trouver la faille.} Un backtest d'une stratégie d'\omterm{def:m1:auctions:call}{enchères} à la clôture suppose que
ses \omterm{def:m1:auctions:moc}{ordres à cours limité à la clôture} (LOC), saisis au \omterm{def:m1:auctions:imbalance}{cours indicatif} affiché à 15:55,
sont exécutés en totalité chaque fois que le \omterm{def:m1:auctions:moc}{cours de clôture} atteint ou dépasse leur limite.
Identifier deux raisons pour lesquelles les exécutions sont surestimées, ainsi que
l'information nécessaire pour les corriger.
\end{exercise}

\section{Problème~: la clôture du rééquilibrage d'indice}

\begin{problem}[{Problème du week-end --- une enchère à la clôture le jour du rééquilibrage}]\label{pb:m1:auctions:1}
Une action entre dans un indice à la clôture du jour. Dix minutes avant la clôture, son
marché continu est à $39.99 \times 40.01$ et la dernière transaction, qui sert de prix de référence,
est à 40.00~; le tick vaut 1 centime. Le carnet de clôture contient, outre
l'ordre des fonds indiciels, les ordres suivants.

Achats~: 50\,000 au marché~; 80\,000 à 40.10~; 120\,000 à 40.00~; 150\,000 à 39.90.
Ventes~: 60\,000 au marché~; 100\,000 à 39.95~; 140\,000 à 40.05~; 200\,000 à
40.15~; 250\,000 à 40.30~; 300\,000 à 40.50.

\medskip\noindent\textbf{Partie I --- Sans les fonds indiciels.}
\begin{enumerate}
\item Tabuler $D(p)$ et $S(p)$ à chaque prix limite, de 39.90 à 40.50.
\item Donner le volume exécutable à chacun, ainsi que le \omterm{def:m1:auctions:call}{prix d'équilibre} et le volume.
\item Donner l'excédent et son côté.
\item Quels ordres sont exécutés partiellement~?
\end{enumerate}

\medskip\noindent\textbf{Partie II --- Les fonds indiciels arrivent.}
Les fonds indiciels ajoutent un ordre d'achat au marché à la clôture de 500\,000 actions.
\begin{enumerate}[resume]
\item Recalculer $D(p)$.
\item Donner le nouveau \omterm{def:m1:auctions:call}{prix d'équilibre} et le volume.
\item De combien, en centimes et en points de base, l'ordre indiciel a-t-il fait monter
le \omterm{def:m1:auctions:moc}{cours de clôture}~?
\item Quel déséquilibre a été publié à 15:50 si l'ordre indiciel était déjà
présent et que le \omterm{def:m1:auctions:call}{prix d'équilibre} avait été calculé au prix de référence~?
\item Combien de dollars les fonds indiciels ont-ils payés au-dessus du prix de la question
2, sur leurs 500\,000 actions~?
\end{enumerate}

\medskip\noindent\textbf{Partie III --- La liquidité arrive.}
Voyant le déséquilibre, les fournisseurs de liquidité ajoutent des ordres de vente~: 150\,000 à 40.10
et 200\,000 à 40.20.
\begin{enumerate}[resume]
\item Recalculer $S(p)$ ainsi que le \omterm{def:m1:auctions:call}{prix d'équilibre} et le volume.
\item De combien les ordres de contrepartie ont-ils fait économiser aux fonds indiciels~?
\item Un fournisseur a vendu 150\,000 actions au nouveau \omterm{def:m1:auctions:moc}{cours de clôture} et s'attend à ce que le titre
revienne à 40.05 le lendemain. Quel est son profit espéré~?
\item Quel risque prend-il pendant la nuit, et comment pourrait-il se couvrir~?
\item Estimer la profondeur de l'\omterm{def:m1:auctions:call}{enchère} $\lambda_D + \lambda_S$ par tick autour du
prix final à partir des tableaux obtenus, et vérifier la \cref{prop:m1:auctions:depth} sur
le mouvement causé par l'ordre indiciel de la partie III.
\end{enumerate}

\medskip\noindent\textbf{Partie IV --- Jugement.}
\begin{enumerate}[resume]
\item Les fonds indiciels auraient pu acheter les jours précédents à la place.
Qu'auraient-ils gagné et qu'auraient-ils risqué~?
\item Qui, économiquement, paie les fournisseurs de liquidité de la question 12~?
\item Un trader saisit un gros ordre d'achat à cours limité à la clôture à 15:57 et
l'annule à 15:59:50 sur une bourse qui l'autorise. Quel était son but, et comment
appelle-t-on cette pratique~?
\item Pourquoi les bourses restreignent-elles la saisie et l'annulation des ordres de clôture
après la première publication du déséquilibre~?
\item Énoncer le \emph{résultat nommé}~: le \omterm{def:m1:auctions:call}{prix d'équilibre} et le volume apparié
de l'\omterm{def:m1:auctions:call}{enchère} finale de la partie III.
\item En une phrase~: qu'est-ce qu'un fonds indiciel achète avec la prime qu'il paie
à la clôture~?
\end{enumerate}
\end{problem}

\section{Questions d'entretien}

\begin{interviewq}[$\star$ \iqroles{trader, développeur, chercheur}]\label{iq:m1:auctions:1}
Comment le prix d'une \omterm{def:m1:auctions:call}{enchère} de clôture est-il déterminé~?
\end{interviewq}

\begin{interviewq}[$\star$ \iqroles{trader, chercheur}]\label{iq:m1:auctions:2}
Pourquoi tant de volume s'échange-t-il à la clôture~?
\end{interviewq}

\begin{interviewq}[$\star\star$ \iqroles{développeur}]\label{iq:m1:auctions:3}
Implémentez un algorithme d'uncrossing dont le temps d'exécution est en $O(n \log n)$ pour $n$ ordres.
Quels sont les cas limites~?
\end{interviewq}

\begin{interviewq}[$\star\star$ \iqroles{trader, chercheur}]\label{iq:m1:auctions:4}
Le déséquilibre de clôture publié pour une action s'élève à 800\,000 actions à l'achat, face
à une \omterm{def:m1:auctions:call}{enchère} moyenne de 1.5 million. Qu'attendez-vous qu'il se produise dans les
dix prochaines minutes, dans l'\omterm{def:m1:auctions:call}{enchère} et sur le marché continu~?
\end{interviewq}

\begin{interviewq}[$\star\star$ \iqroles{chercheur, ingénieur ML}]\label{iq:m1:auctions:5}
Vous voulez prédire le \omterm{def:m1:auctions:moc}{cours de clôture} à partir des messages de déséquilibre. Quelle est votre
cible, quelles sont vos variables explicatives et quel est le principal piège de
l'évaluation~?
\end{interviewq}

\begin{interviewq}[$\star\star\star$ \iqroles{trader, chercheur}]\label{iq:m1:auctions:6}
Une plateforme propose de remplacer la négociation en continu par une \omterm{def:m1:auctions:call}{enchère} toutes les 100
millisecondes. Qui y gagne, qui y perd, et qu'attendez-vous qu'il advienne des
fourchettes~?
\end{interviewq}
